\chapter{Optimization}
\label{sec:Optimization}

The following algorithms are used to minimize a function.

\subsection{Grid} \label{sec:GridSearch}

A grid search has a number of dimensions, called N. Each dimension contains a number of numeric items.
All item combinations, also called points, are evaluated with the cost function.
In essence, the grid search does no more that sequentially evaluating all these points and return the point with the minimum cost value.

Each dimension has a number of items which should be evaluated.
The items are defined as a range with a minimum and maximum value and the number of steps into which the range should be subdivided.

Possibly one item cannot be evaluated, because the realization can not be calculated.
It will be ignored in the optimization process.

\subsubsection{Move grid if minimum is on edge} \label{sec:MoveGrid}

If, in any dimension, the optimum is on the edge of a dimension,
the edge will be extended one step in the direction of the edge.
This extension can happen in several dimensions at the same time. 

Only dimensions with at least three items can be moved and the items should be defined as minimum/maximum/steps.
Per dimension this feature is optional.

This process will be repeated until the optimum is not on the edge of the dimension area any more.

\subsubsection{Refinement of the grid}\label{sec:RefineGrid}

When the minimum has been found, eventually moved (\autoref{sec:MoveGrid}),
the optimization process will refine the grid automatically around the point with the optimum.
Per dimension three items will be defined, around the point and with half of the step size.
This leads to a $3^{N}-1$ new evaluations in an N-dimensional grid. 

Only dimensions with at least two items can be refined and the items should be defined as minimum/maximum/steps.
Per dimension this feature is optional.

Several refinement steps can be taken. Refinement stops after a provided number of refinement steps.

\begin{figure}[H]
	\centering
	\includegraphics[width=0.8\textwidth]{\pathProbLibPictures RefinementGrid.png}
	\caption{Refinement grid in a two dimensional grid}
	\label{fig:RefinementGrid}
\end{figure}

\subsection{Genetic algorithm} \label{sec:GeneticAlgorithm}

The genetic algorithm is a population-based, derivative-free optimization procedure.
It is useful for objective functions that are discontinuous, non-differentiable, stochastic, strongly non-linear, or dependent on a large number of variables.
The procedure followed by the implementation is shown in \autoref{fig:Figure_ProcedureGA}.

A population consists of candidate parameter vectors, called individuals.
Each individual is evaluated by the objective function, and the objective value is used as the individual's fitness.
The implementation minimizes this fitness value, so a lower value represents a better individual.

The population evolves in generations.
The initial population is generated by the selected initialization method.
At the start of each generation, the best individuals are retained as elites.
The remaining population is replaced by children generated from two randomly selected parents.
Each child is formed through crossover and can be modified by mutation.
Two children are generated for every non-elite position, and the child with the lower objective value is retained.

The process continues for the configured maximum number of generations.
It also stops after a configured number of generations without improvement or earlier if the model runner requests termination.
Elitism ensures that the best solution found is not lost while the evolution of the remaining individuals continues.

\begin{figure}[H]
	\centering
	\includegraphics[width=0.30\textwidth]{\pathProbLibPictures GeneticAlgorithm.png}
	\caption{Procedure of the genetic algorithm}
	\label{fig:Figure_ProcedureGA}
\end{figure}

\subsubsection{Initialization}  
\label{sec:GeneticAlgorithmInitialization}
Like in any other optimization method an initial guess of the final outcome is required.
This guess concerns a choice of all individual genes in the population.
In many problems the influence of the individual genes is fully independent.
By implication, a random choice of the values of the genes between their appropriate limits is adequate.

In some problems a weak connection between the gene values is observed.
An example is a free slip plane in slope analysis.
Such a plane must be as flexible as possible in order to allow a wide variety of shapes.
However, buckling planes will not survive as they have a surplus of resisting power.
It is better to exclude these already at initialization.

From these perspectives, so far, two ways of initialization are defined:
\begin{itemize}
	\item Uniform:
	\begin{align}
	\frac{g_{n} - g_{min}}{g_{max} - g_{min}} = \eta_{n}  \quad \text{with} \quad 0 < \eta_{n} < 1
	\end{align}
	\item Parabolic:
	\begin{align}
	\nonumber &\dfrac{g_{n} - g_{min}}{g_{max} - g_{min}} = \\
	&\left\{ \begin{array}{lll}
	\left(\eta_{A} - \eta_{B}\right) \left(\dfrac{1}{\eta_{C}} \dfrac{n-1}{N-1} - 1 \right)^{2} + \eta_{B} & \text{if} & \dfrac{1}{2} < \eta_{C} < \infty \\
	\left(\eta_{A} - \eta_{B}\right) \left(\dfrac{1}{1 - \eta_{C}} \dfrac{N - n}{N-1} - 1 \right)^{2} + \eta_{B} &  \text{if} & -\infty < \eta_{C} < \dfrac{1}{2} 	\end{array} \right. \\
	\nonumber &\text{with} \quad \quad 0 < \eta_{A} < 1 \quad \text{and} \quad 0 < \eta_{B} < 1
	\end{align}
\end{itemize}

In the uniform case every gene has its own random number.
The corresponding value of the gene ranges between its minimum and maximum value.
These bounds may even be different for every gene.
The parabolic case is defined by three random numbers.
They are chosen in such a way that the parabola for all genes ranges between 0 and 1.

Note that $\eta_{c}$ is the relative position of the parabola's extreme.
If this position is located outside the gene's position, the available space is not uniformly filled up, since the extreme enforces an extra offset.
Therefore, it could be decided to fix the extreme always in the gene's range, $0 < \eta_{c} < 1$.

\subsubsection{Parental selection} 
\label{sec:ParentalSelection}
Selection specifies how the genetic algorithm chooses parents for the next generation.
The implemented selection method is a uniform random choice within the population.
Two parent indices are selected independently:
\begin{align}
P = G_{\text{int}\left(\eta M\right)},
\end{align}
where $\eta$ is a random number between zero and one, $M$ is the population size, and $G$ is the current population.

No fitness-based scaling or weighted selection is applied.
The same population can therefore contain both well-performing and poorly performing individuals as parents.

\subsubsection{Crossover}

The crossover rules combine two parents to form a child.
The crossover mask is a vector of choices: $0$ selects the first randomly selected parent and $1$ selects the second.
The implementation provides the following methods:

\begin{itemize}
	\item Scattered:
	
	The father and mother are combined by having a completely random crossover mask.
	This works well for a small genome with independent genes. The crossover mask reads:
	\begin{align}
	\chi =\left( 0 | 1 \right)_{n}
	\end{align}
	
	\item Single Point:
	A single point crossover will combine the first part of one parent with the second part of the other parent.
	For genome types which specify a coherent object it will facilitate a more smooth convergence. It is defined as:
	\begin{align}
	\chi = \left\{0 ... 0, 1 ... 1\right\}
	\end{align}
	
	\item Double point:
	A double point crossover combines the middle of one slip plane with the edges of another. The definition is:
	\begin{align}
	\chi = \left\{0 ... 0, 1 ... 1, 0 ... 0\right\}
	\end{align}
	
\end{itemize}

\subsubsection{Mutation} 
\label{sec:Mutation}

Mutation makes random changes in a child and provides genetic diversity.
The implementation defines three mutation methods:

\begin{itemize}
	\item Jump: sets each selected gene to a uniformly random value within its bounds.
	A jump helps finding the global minimum by jumping out of a local minimum. 
	\item Creep: changes a gene relative to the mutation anchor, by at most the configured creep reduction times the gene range.
	\item Inversion: reflects a gene across the lower and upper bounds.
\end{itemize}

Each mutation method may be assigned a fraction that determines how often it is selected.
The crossover and mutation fractions can be randomized separately for every generation or every child.
After crossover and mutation, the child is checked against fixed lower and upper bounds.

\subsubsection{Child selection} \label{sec:ChildSelection}

For each non-elite position, the routine creates two children and retains the child with the lower objective value.
This local tournament between paired children favors improvement while preserving diversity from the remaining population.

The configured elite is passed directly into the next generation.
It is kept at the beginning of the population and ordered according to objective value.
The elite size and population size are configurable.

The search terminates after the maximum number of generations or after a configured number of generations without improvement.
It can also stop earlier when the model runner requests termination.

\subsection{Adaptive Particle Swarm Optimization} 
\label{sec:APSO}

Adaptive Particle Swarm Optimization (APSO) is a population-based optimization method that combines an APSO search with differential evolution.
Each particle represents a candidate vector of optimization parameters.
The parameter values are considered in normalized form, with values between zero and one, while model evaluations use the corresponding absolute parameter values.
The implementation starts with a population generated by the parabolic initialization method and evaluates every particle.
The particle with the lowest objective value is stored as the best particle found so far.

At every generation, the particles are ranked by objective value and the best configured number of particles is retained as the elite.
New positions are calculated from an APSO branch or a differential-evolution branch.
The probability of selecting the APSO branch decreases from approximately $0.9$ to $0.1$ over the course of the search.

The step from one population to the next is given by

\begin{align}\label{eq:APSO}
X_{i\text{,APSO}}^{t+1} = \left\{\begin{array}{ll}r_{\text{c}} < C & \left(1 - \beta\right)X_{i}^{t} + \beta B_{i}^{t} + \alpha r_{\text{a}} \\
r_{\text{c}} > C & X_{i}^{t}
\end{array}	
\right.
\end{align}

with

\begin{align}
\alpha = \delta ^ t
\end{align}

where

\begin{longtable*}{p{18mm}p{\textwidth-24pt-18mm}}  
	$i$ & is the index number of the parameter in the particle; \\  
	$X$ & represents the particle; \\  
	$t$ & is the generation step in the iteration process;\\
	$B$ & is the best particle found so far;\\  
	$C$ & is a given cross over fraction [0, 1];\\  
	$\beta$ & is the parameter convergence factor [0, 1]; \\  
	$\delta$ & is the parameter step size [0.1, 0.99]; \\  
	$r_{\text{c}}$ & is a random number, generated per parameter, particle and generation [0, 1]; \\
	$r_{\text{a}}$ & is a random number, generated per parameter, particle and generation [-1, 1]; 
\end{longtable*}

The algorithm keeps an elite, which consists of the best particles found so far.
The elite has size $K$.
To improve the population,
the APSO algorithm is combined with a differential evolution algorithm, which is defined by

\begin{align}\label{eq:DE}
X_{\text{DE,} i}^{t+1} = \left\{\begin{array}{ll}r_{\text{c}} < C & B_{i}^{t} + F \times \left(E_{\text{r1,} i}^{t} - E_{\text{r2,} i}^{t} + E_{\text{r3,} i}^{t} - E_{\text{r4,} i}^{t}\right)\\
r_{\text{c}} > C & X_{i}^{t}
\end{array}	
\right.
\end{align}

where

\begin{longtable*}{p{18mm}p{\textwidth-24pt-18mm}}  
	$F$ & is the given differential weight; \\  
	$E$ & are particles in the elite; \\  
	$r1, r2, r3, r4$ & are random index numbers in the elite; \\  
\end{longtable*}

Combining these algorithms, the next particle is defined as 

\begin{align}\label{eq:APSODE}
X_i^{t+1} = \left\{\begin{array}{ll}r_{\text{s}} < S & X_{\text{APSO,} i}^{t+1}\\
r_{\text{s}} > S & X_{\text{DE,} i}^{t+1}
\end{array}	
\right.
\end{align}

where

\begin{longtable*}{p{18mm}p{\textwidth-24pt-18mm}}  
	$S$ & is the APSO/DE separation factor, which linearly declines from 0.9 in the first generation
	to 0.1 in the last generation; \\  
	$r_{\text{s}}$ & is a random number, generated per particle and generation [0, 1]; \\  
\end{longtable*}

Using a random factor, in each iteration and per each particle it is decided
whether the new particle comes from \autoref{eq:APSO} or from \autoref{eq:DE}.
This is a random process which uses the APSO/DE separation factor $S$.
Next, another random process uses the crossover fraction $C$ to decide for each parameter whether that value is updated.
Parameters not selected by the crossover process remain unchanged.

The normalized candidate value is converted back to an absolute parameter value before model evaluation.
If a proposed value lies outside a fixed lower or upper bound, the particle uses the current best-particle value for that parameter instead.

After evaluating a generation, the implementation updates the best particle and reports convergence.
The search stops after the configured maximum number of generations or after a configured number of generations without improvement.
It can also stop earlier if the model runner requests termination.

\subsection{Cobyla}
\label{sec:CobylaOptimization}

Cobyla (Constrained Optimization by Linear Approximations) is a derivative-free optimization algorithm for problems with a nonlinear objective function and nonlinear constraints.
It is implemented as a sequential trust-region method \cite{Powell1994}.
The algorithm first defines a simplex consisting of the starting point and $n$ additional points for an optimization problem with $n$ variables.
At every iteration, linear approximations of the objective function and the constraint functions are constructed by interpolation through the simplex points.

The interpolated linear model is then used to determine a trial point within a trust region around the best point found so far.
The trial point is accepted if it improves the objective function sufficiently and satisfies the constraints.
Otherwise, the simplex and trust-region radius are adjusted and the iteration is repeated.
This process avoids numerical derivatives and therefore allows Cobyla to handle nonsmooth or discontinuous model responses.

The search is limited by the lower and upper bounds of each optimization variable.
Constrained Cobyla problems require the constraint value to satisfy the constraint tolerance specified in the settings.
The optimization stops when the relative change in the variables becomes sufficiently small or when the maximum number of model evaluations is reached.

\subsection{K-Means clustering}
\label{sec:KMeansClustering}

K-Means clustering partitions a set of samples into $k$ groups such that samples in the same group are close to their cluster center.
The cluster center is the arithmetic mean of the sample values in the cluster.
For a sample $\mathbf{x}_i$ and cluster $j$, the squared Euclidean distance is

\begin{align}
d^2\left(\mathbf{x}_i,\mathbf{c}_j\right)
= \sum_{n=1}^{N}\left(x_{i,n}-c_{j,n}\right)^2 ,
\end{align}

where $\mathbf{c}_j$ is the center of cluster $j$ and $N$ is the number of variables.
Sample weights can be included when the initialization selects the next cluster center.

The implementation initializes the clusters with the K-Means++ principle \cite{Arthur2007}.
The first center is the mean of all samples.
Each subsequent center is selected randomly from the samples, with a probability proportional to its weighted squared distance from the nearest existing center.
This favors center locations in regions that are not yet represented well.

The algorithm then repeatedly assigns every sample to the cluster with the nearest center.
After the assignments, all cluster centers are updated to the mean of their assigned samples.
Iterations continue until the assignments no longer change or the maximum number of iterations is reached.
Multiple trials with different initializations can be performed, and the trial with the smallest sum of squared distances between samples and their cluster centers is retained.

The number of clusters can be fixed by the user or optimized automatically.
In the automatic case, the algorithm evaluates solutions with an increasing number of clusters and uses the silhouette coefficient \cite{Rousseeuw1987} to measure cluster quality.
The silhouette coefficient compares the mean distance within a cluster with the mean distance to the nearest other cluster.
Values range from $-1$ to $1$, with values near $1$ indicating well-separated clusters.
